% !TeX root = ../../relatorio-projeto-exemplo-undf.tex
\chapter{Introdução e contexto}
Este relatório didático descreve a organização dos modelos de relatório final
do UnDF-TeX. A demanda foi substituir um roteiro de planejamento de projeto
por uma estrutura que permitisse relatar desenvolvimento realizado, entregas,
avaliação e limitações. O recorte abrangeu as modalidades de projeto técnico
ou profissional e de modelo de negócio de base tecnológica descritas no
manual em elaboração \citeind{manualundf}.

Para demonstrar a modalidade profissional, o caso adota um cenário fictício
de atendimento documental a uma organização. O estudante é associado à
organização dos arquivos e à verificação das entregas. Não houve contratante,
estágio ou implantação reais; a evidência verificável é a produção documental
local. O recorte preserva a etapa de dois modelos e dois exemplos descrita
neste caso, anterior à inclusão do modelo de Projeto Aplicado.

O objetivo foi disponibilizar entradas compiláveis, conteúdo para preenchimento
e exemplos separados. A contribuição demonstrada é documental e técnica:
não houve pesquisa com estudantes, implantação institucional ou mensuração
de desempenho acadêmico. Os nomes da autoria e da banca são fictícios.
A organização do relato também foi cotejada com a Parte I do relatório
final do UnDF-TeX \citeind{relatoriofinal}, tomada apenas como referência.

\chapter{Fundamentação e requisitos}
A separação entre conteúdo e apresentação foi adotada para permitir que o
estudante editasse seu texto sem alterar a classe. O LaTeX oferece recursos
para estruturação, referências cruzadas e bibliografia, conforme a apresentação
do projeto \citeind{latexproject}. Esses recursos fundamentam a escolha da
ferramenta, mas não demonstram, por si, redução de tempo de escrita.

Foram definidos quatro critérios de aceitação: disponibilizar um modelo e
um exemplo para cada modalidade; manter metadados, resumos e bibliografia
independentes entre trabalhos; preservar os elementos institucionais; e
gerar PDFs com citações e sumário resolvidos. A organização do miolo seguiu
o roteiro de relatório final, sem transformar o texto em plano de atividades.

\chapter{Metodologia e desenvolvimento realizado}
O trabalho consistiu em inspeção dos arquivos existentes, reorganização do
conteúdo, escrita dos roteiros e dos exemplos e compilação local. O antigo
conteúdo de projeto foi transferido para a pasta de modelos e reescrito em
sete seções primárias. A modalidade de negócio tecnológico foi organizada
em oito seções primárias.

A apresentação institucional foi centralizada em um arquivo comum. O
conteúdo editável permaneceu separado por modalidade e por finalidade:
modelo para preenchimento ou demonstração. A classe e os mecanismos de
citações e referências foram reutilizados. Os documentos foram compilados
com LuaLaTeX e Biber, com passagens adicionais para resolver o sumário.

\section{Comparação entre previsto e realizado}
O balanço da execução distinguiu as entregas concluídas dos limites do recorte.
A Tabela~\ref{tab:balanco-profissional} registra essa relação sem atribuir
uma validação em ambiente profissional aos testes documentais.

\begin{undftabela}[H]{Balanço da execução no caso didático}{balanco-profissional}[Elaboração própria a partir dos arquivos e verificações locais]
\begin{tabular}{@{}p{0.28\linewidth}p{0.29\linewidth}p{0.33\linewidth}@{}}
\hline
\textbf{Item} & \textbf{Realizado} & \textbf{Evidência ou limite} \\
\hline
Modelos e exemplos & Quatro entradas produzidas & Fontes e PDFs compilados \\
Apresentação institucional & Arquivo comum reutilizado & Composição dos quatro documentos \\
Uso por organização & Fora do recorte & Ausência de implantação e pesquisa com pessoas \\
\hline
\end{tabular}
\end{undftabela}

\chapter{Solução e entregas}
Foram produzidas quatro entradas: dois modelos para preenchimento e dois
exemplos didáticos. A entrada \texttt{relatorio-projeto.tex} corresponde ao
relatório técnico ou profissional e \texttt{negocio-tecnologico.tex}, ao
relatório de desenvolvimento e análise do modelo de negócio. Cada entrada
possui um exemplo identificado pelo sufixo \texttt{-exemplo-undf}.

Os arquivos de cada trabalho foram separados em metadados, resumos, miolo
e bibliografia. Um índice explica a escolha da entrada, os caminhos de
edição e a compilação. Os roteiros preservam a possibilidade de adaptar
subdivisões conforme o objeto de pesquisa e a orientação do curso.

\chapter{Resultados e avaliação}
A avaliação local verificou a compilação das quatro entradas e a presença
de resumo, abstract, sumário e referências. Também conferiu a organização
em sete seções primárias no projeto e oito no negócio tecnológico, além
da abertura das referências em página própria. Esses critérios avaliam
o funcionamento documental, sem substituir a revisão acadêmica do conteúdo.

O resultado foi um conjunto de modelos preenchíveis acompanhado de
demonstrações independentes. A revisão do conteúdo distinguiu entregas,
testes e efeitos observados. No exemplo de negócio, os valores financeiros
foram identificados como premissas de simulação, sem alegação de coleta
com clientes. Logs e PDFs gerados constituem evidências locais reproduzíveis.

\chapter{Discussão e limitações}
A reutilização da apresentação reduziu a duplicação de comandos institucionais.
Em contrapartida, uma alteração nesse arquivo comum afeta os quatro
documentos e exige nova verificação. O compartilhamento foi limitado à
apresentação: os textos e as referências permanecem independentes.

Não foram conduzidos testes de usabilidade nem avaliação por banca ou
biblioteca. A adequação às exigências específicas de cada curso permanece
dependente de conferência institucional. Os exemplos são curtos e didáticos;
sua extensão não define a extensão adequada de um TCC real.

\chapter{Conclusão}
O objetivo técnico foi atendido com dois modelos de relatório final e seus
exemplos, organizados para edição e compilação independentes. O roteiro
anterior de planejamento foi substituído por uma estrutura orientada ao
relato do trabalho realizado e à discussão de seus resultados.

A contribuição está na organização dos artefatos e na coerência entre o
manual e os modelos. Recomenda-se, como continuidade, revisar as exigências
por curso e observar o uso dos modelos por estudantes. Essas atividades
constituem recomendações; não foram realizadas nem usadas para afirmar
eficácia institucional neste relatório.
